\section{引言}

数据要素流通要求数据在跨机构、跨平台和跨系统使用过程中保持可信。这里的可信不只表示底层文件没有被篡改，还包括授权过程可追溯、计算任务可复核、审计证据不可抵赖，以及异常发生后能够界定责任。金融、医疗、政务和工业互联网等场景通常不允许数据明文自由流动，而是倾向于采用隐私计算、可信执行、数据连接器或受控接口完成协作。因此，数据流通中的校验对象已经从静态文件扩展为数据分片、计算日志、模型更新、任务摘要和结果哈希等多类证据。

传统云存储审计主要回答“远程数据是否仍被正确保存”的问题。可证明数据持有（Provable Data Possession, PDP）通过挑战响应和同态标签降低完整性验证成本~\cite{ateniese2007pdp}；可恢复性证明（Proofs of Retrievability, PoR）进一步强调数据在遭受部分损坏后仍应可恢复~\cite{juels2007por,shacham2013por}；动态 PDP 支持数据块插入、删除和修改，使完整性证明能够覆盖动态外包数据~\cite{erway2015dpdp}。这些工作构成了远程完整性校验的基础，但其审计过程通常仍依赖第三方审计员或中心化服务，审计行为本身缺少不可篡改记录。

区块链被引入数据审计后，研究重点从“数据能否被验证”扩展为“审计过程是否可信”。已有工作使用区块哈希或随机数生成不可预测挑战，利用智能合约记录审计结果，或通过联盟链共识约束审计员拖延和合谋行为~\cite{xue2019ibpa,zhang2021cpvpa,fan2020dredas,lin2022cbpiv}。在这一脉络中，BLoP 进一步把区块链、可信计算和隐私保护计算整合起来，通过可信度量、可信 Agent、异常事件监测和 PBFT 风格共识，形成面向可信数据流通的全节点审计模式~\cite{blop2024trusted}。这种全节点模式适合对安全性和一致性要求较高的场景，可以作为本文轻量化设计的基础模式和对照基线。

本文围绕“基于区块链的分布式数据校验方法”提出一个小规模、可验证的轻量化升级设计：TrustFlowAudit。该方法继承 BLoP 的可信节点治理和 PBFT 链上确认思路，但不要求每一次审计都进入全节点确认流程，而是将 BLoP 全节点模式视为高安全等级基线，在中低风险或高频审计场景中提供可信委员会、批量存证和纠删码恢复组成的轻量化路径。本文聚焦一个可闭环的问题：当链下数据分片或审计证据被篡改、丢失或异常提交时，系统如何低成本发现异常、触发恢复、记录责任并更新节点信誉。本文的核心思想是：链下完成数据分片、哈希校验、纠删码恢复和证据保存，链上只保存必要承诺、状态摘要和责任记录；同时基于可信评分选择小规模审计委员会，作为 BLoP 全节点确认模式之外的可选升级策略。

本文主要贡献如下：

\begin{itemize}
  \item 在 BLoP 全节点可信审计模式基础上，提出一个面向高频审计场景的轻量化升级路径，将可信节点集合、PBFT 确认和审计证据存证扩展为可分级执行的流程。
  \item 提出面向数据流通场景的分布式数据校验闭环，将 PDP/PoR 思想、区块链存证、纠删码恢复和联盟链确认统一到“发现、恢复、记录、惩罚”的流程中。
  \item 设计可信评分驱动的轻量审计委员会机制，在保留 BLoP 全节点模式作为高安全基线的同时，为中低风险审计提供小委员会确认选项。
  \item 给出链下计算与链上存证的分层架构：链下保存原始证据和执行复杂校验，链上记录 Merkle 根、哈希链摘要、审计状态和信誉变化。
  \item 明确方法边界和后续实验路线，为后续在 BLoP 机制复现、联盟链容器化验证和纠删码审计恢复实验中开展定量评估提供基础。
\end{itemize}
